\documentclass{article}


\usepackage{PRIMEarxiv}

\usepackage[utf8]{inputenc} % allow utf-8 input
\usepackage[T1]{fontenc}    % use 8-bit T1 fonts
\usepackage{hyperref}       % hyperlinks
\usepackage{url}            % simple URL typesetting
\usepackage{booktabs}       % professional-quality tables
\usepackage{amsfonts}       % blackboard math symbols
\usepackage{nicefrac}       % compact symbols for 1/2, etc.
\usepackage{microtype}      % microtypography
\usepackage{lipsum}
\usepackage{fancyhdr}       % header
\usepackage{graphicx}       % graphics
\graphicspath{{media/}}     % organize your images and other figures under media/ folder

%Header
\pagestyle{fancy}
\thispagestyle{empty}
\rhead{ \textit{ }} 

% Update your Headers here
\fancyhead[LO]{Running Title for Header}
% \fancyhead[RE]{Firstauthor and Secondauthor} % Firstauthor et al. if more than 2 - must use \documentclass[twoside]{article}



  
%% Title
\title{Substitution Artifacts in Shapley-Based Temporal Attribution:
Characterizing and Correcting Masking-Background Bias on
Mean-Nonstationary Time Series
%%%% Cite as
%%%% Update your official citation here when published 
\thanks{\textit{\underline{Citation}}: 
\textbf{Authors. Title. Pages.... DOI:000000/11111.}} 
}

\author{
  Myroslav Mishchuk \\
  Silesian University of Technology \\ 
  Gliwice, Poland \\
  Lviv Polytechnic National University \\ 
  Lviv, Ukraine \\
  \texttt{myroslav.mishchuk@polsl.pl} \\
   \And
  Rafał Cupek \\
  Silesian University of Technology \\ 
  Gliwice, Poland\\
  \texttt{rafal.cupek@polsl.pl} \\
    \And
  Olena Pavliuk \\
  Silesian University of Technology \\ 
  Gliwice, Poland \\
  Lviv Polytechnic National University \\ 
  Lviv, Ukraine \\
  \texttt{olena.pavliuk@polsl.pl} \\
}


\begin{document}
\maketitle


\begin{abstract}
\lipsum[1]
\end{abstract}


% keywords can be removed
\keywords{explainable AI \and Shapley values \and temporal attribution \and time series \and mean-nonstationarity \and remaining useful life}


\section{Introduction}
\label{sec:intro}

\subsection{Research questions}\label{sec:setup:rqs}
\begin{itemize}
    \item \textbf{RQ1}: How large and how structured is the distortion that default backgrounds (zeros, global mean) induce in temporal Shapley attributions on mean-nonstationary data?
    \item \textbf{RQ2}: Do trend-aware backgrounds reduce the distortion, and how much of the gap to the learned-imputer ceiling do they close? 
    \item \textbf{RQ3}: Do the synthetic-benchmark verdicts transfer to the real-world datasets?
\end{itemize}

\section{Related Work}
\label{sec:related}

\subsection{The value-function problem in Shapley explanations}
\label{sec:related:offmanifold}

In the Shapley framework, the value of a feature coalition is computed by evaluating the model with the complementary features replaced or marginalized, and the choice of this value function is known to change the resulting explanation. \citet{frye2021shapley} demonstrated that ``off-manifold'' evaluation --- splicing replacement values into an observed instance without regard for feature dependence --- can produce incorrect explanations, and proposed two on-manifold remedies: substitution via a generative model of the data, and directly learning the value function. \citet{aas2021explaining} developed conditional-expectation approximations for dependent features in the kernel-SHAP setting. The conceptual axis along which these proposals differ is now standard: interventional value functions stay true to the model but evaluate it on unrealistic inputs, while observational/conditional value functions respect the data distribution at the cost of attributing importance to merely correlated features \citep{chen2020true, janzing2020feature, sundararajan2020many}. \citet{yeh2022threading} formalized desiderata spanning both views and showed that both classes are vulnerable to manipulation in low-density regions. The line remains active: recent work casts on-manifold attribution as transport along generative flows \citep{zhang2026axiomatic}.
 
Two properties of this literature matter for the present paper. First, its data model is static and tabular (or, in the recent theoretical work, image-like): the manifold is generic, and the proposed remedies therefore require estimating conditional distributions or training generative surrogates. Second, none of it examines temporally ordered data in which the manifold possesses a \emph{known, low-dimensional organizing variable} --- the position a sequence has reached along its mean trajectory --- which, as we argue, simultaneously structures the off-manifold bias and makes conditioning cheap.
 

\subsection{Off-manifold Shapley on tabular data}\label{sec:related:tabular}

\subsection{Masking artifacts and realistic infilling in image XAI}\label{sec:related:image}

\subsection{Realistic perturbations for time-series saliency}\label{sec:related:ts}

\subsection{Temporal Shapley methods and background sensitivity}\label{sec:related:tshap}

\subsection{XAI in prognostics}\label{sec:related:phm}

\subsection{The gap}\label{sec:related:gap}

\lipsum[4] See Section \ref{sec:headings}.

\subsection{Headings: second level}
\lipsum[5]
\begin{equation}
\xi _{ij}(t)=P(x_{t}=i,x_{t+1}=j|y,v,w;\theta)= {\frac {\alpha _{i}(t)a^{w_t}_{ij}\beta _{j}(t+1)b^{v_{t+1}}_{j}(y_{t+1})}{\sum _{i=1}^{N} \sum _{j=1}^{N} \alpha _{i}(t)a^{w_t}_{ij}\beta _{j}(t+1)b^{v_{t+1}}_{j}(y_{t+1})}}
\end{equation}

\subsubsection{Headings: third level}
\lipsum[6]

\paragraph{Paragraph}
\lipsum[7]

\section{Examples of citations, figures, tables, references}
\label{sec:others}
\lipsum[8] \cite{kour2014real,kour2014fast} and see \cite{hadash2018estimate}.

The documentation for \verb+natbib+ may be found at
\begin{center}
  \url{http://mirrors.ctan.org/macros/latex/contrib/natbib/natnotes.pdf}
\end{center}
Of note is the command \verb+\citet+, which produces citations
appropriate for use in inline text.  For example,
\begin{verbatim}
   \citet{hasselmo} investigated\dots
\end{verbatim}
produces
\begin{quote}
  Hasselmo, et al.\ (1995) investigated\dots
\end{quote}

\begin{center}
  \url{https://www.ctan.org/pkg/booktabs}
\end{center}


\subsection{Figures}
\lipsum[10] 
See Figure \ref{fig:fig1}. Here is how you add footnotes. \footnote{Sample of the first footnote.}
\lipsum[11] 

\begin{figure}
  \centering
  \fbox{\rule[-.5cm]{4cm}{4cm} \rule[-.5cm]{4cm}{0cm}}
  \caption{Sample figure caption.}
  \label{fig:fig1}
\end{figure}

\subsection{Tables}
\lipsum[12]
See awesome Table~\ref{tab:table}.

\begin{table}
 \caption{Sample table title}
  \centering
  \begin{tabular}{lll}
    \toprule
    \multicolumn{2}{c}{Part}                   \\
    \cmidrule(r){1-2}
    Name     & Description     & Size ($\mu$m) \\
    \midrule
    Dendrite & Input terminal  & $\sim$100     \\
    Axon     & Output terminal & $\sim$10      \\
    Soma     & Cell body       & up to $10^6$  \\
    \bottomrule
  \end{tabular}
  \label{tab:table}
\end{table}

\subsection{Lists}
\begin{itemize}
\item Lorem ipsum dolor sit amet
\item consectetur adipiscing elit. 
\item Aliquam dignissim blandit est, in dictum tortor gravida eget. In ac rutrum magna.
\end{itemize}


\section{Conclusion}
Your conclusion here

\section*{Acknowledgments}
This was was supported in part by......

%Bibliography
\bibliographystyle{unsrt}  
\bibliography{references}  


\end{document}
